\section{ТЕОРЕТИЧНІ ОСНОВИ ТА УМОВИ МОДЕЛЮВАННЯ}

\paragraphheading{Постановка задачі.}
Необхідно дослідити виконання розподіленої транзакції на прикладі банківського переказу між рахунками, які розміщені на різних вузлах. У практичній моделі \texttt{Node A} зберігає рахунок відправника, \texttt{Node B} --- рахунок отримувача, а координатор приймає єдине глобальне рішення: \texttt{COMMIT} або \texttt{ROLLBACK}. Часткове виконання операції є неприпустимим, оскільки списання коштів без відповідного зарахування порушує цілісність даних.

Два вузли моделюються окремими базами даних PostgreSQL: \texttt{bank\_node\_a} та \texttt{bank\_node\_b}. Координатор реалізовано як Python-програму, що встановлює підключення до обох баз і використовує команди \texttt{PREPARE TRANSACTION}, \texttt{COMMIT PREPARED} та \texttt{ROLLBACK PREPARED}. У реальній розподіленій системі ці бази можуть функціонувати на різних серверах.

Для використання prepared transactions параметр PostgreSQL \texttt{max\_prepared\_transactions} має бути більшим за нуль. Підготовлені транзакції не потрібно залишати незавершеними протягом тривалого часу, оскільки вони можуть утримувати блокування та зменшувати доступність ресурсів.

\paragraphheading{Порівняння протоколів фіксації.}
Основні характеристики підходів до завершення розподілених транзакцій наведено в табл.~\ref{tab:protocols}.

\begin{table}[htbp]
\centering
\caption{Порівняльний аналіз протоколів фіксації транзакцій}
\label{tab:protocols}
\fontsize{10pt}{12pt}\selectfont
\setlength{\tabcolsep}{3pt}
\renewcommand{\arraystretch}{1.16}
\begin{tabularx}{\textwidth}{|>{\RaggedRight\arraybackslash}p{0.15\textwidth}|
                               Y|Y|Y|
                               >{\RaggedRight\arraybackslash}p{0.14\textwidth}|}
\hline
\textbf{Протокол} & \textbf{Принцип роботи} & \textbf{Переваги} & \textbf{Недоліки} & \textbf{Рівень надійності} \\
\hline
Двофазний протокол (2PC) &
Фаза \texttt{PREPARE}: усі вузли перевіряють операцію та голосують \texttt{READY}/\texttt{ABORT}. Після цього координатор надсилає \texttt{COMMIT} або \texttt{ROLLBACK}. &
Забезпечує атомарне спільне рішення; підтримується PostgreSQL. &
Після \texttt{READY} учасник утримує ресурси до рішення координатора; збій координатора знижує доступність. &
Високий щодо атомарності; залежить від відновлення координатора. \\
\hline
Трифазний протокол (3PC) &
Містить проміжну фазу \texttt{PRE-COMMIT} між голосуванням та остаточною фіксацією. &
Зменшує ризик тривалого блокування за визначених припущень щодо тайм-аутів і мережі. &
Потребує більшої кількості повідомлень; не усуває всі проблеми поділу мережі. &
Високий у визначеній моделі відмов. \\
\hline
Оптимістичні методи &
Операції виконуються без тривалих блокувань, а перед фіксацією виконується перевірка конфліктів. &
Забезпечують добру продуктивність за малої кількості конфліктів. &
За конкуренції спричиняють скасування та повторні спроби; розподілена валідація є складнішою. &
Залежить від алгоритму валідації та повторення операцій. \\
\hline
\end{tabularx}
\end{table}

\paragraphheading{Властивості транзакцій у розподіленому середовищі.}
Застосування 2PC пов'язане із забезпеченням властивостей ACID. Способи їх підтримання в обраній моделі подано в табл.~\ref{tab:acid}.

\begin{table}[htbp]
\centering
\caption{Забезпечення властивостей ACID у розподіленій системі}
\label{tab:acid}
\fontsize{11pt}{13pt}\selectfont
\setlength{\tabcolsep}{4pt}
\renewcommand{\arraystretch}{1.18}
\begin{tabularx}{\textwidth}{|>{\RaggedRight\arraybackslash}p{0.18\textwidth}|
                               >{\RaggedRight\arraybackslash}p{0.29\textwidth}|
                               Y|}
\hline
\textbf{Властивість} & \textbf{Зміст} & \textbf{Спосіб забезпечення у моделі} \\
\hline
Атомарність &
Зміни виконуються на всіх вузлах або не виконуються ніде. &
Координатор 2PC; \texttt{PREPARE TRANSACTION} на кожному вузлі; єдине рішення \texttt{COMMIT PREPARED} або \texttt{ROLLBACK PREPARED}. \\
\hline
Узгодженість &
Після завершення транзакції дані відповідають правилам предметної області. &
Обмеження \texttt{CHECK}, \texttt{PRIMARY KEY}, перевірка активності рахунку та достатності коштів до стану \texttt{READY}. \\
\hline
Ізольованість &
Паралельні операції не повинні формувати некоректний результат. &
Рівень \texttt{SERIALIZABLE}, блокування змінюваних рядків і механізм MVCC PostgreSQL. \\
\hline
Довговічність &
Зафіксовані зміни зберігаються після збою. &
Журнал WAL та збережений стан prepared transaction; після відновлення координатор завершує зафіксоване рішення. \\
\hline
\end{tabularx}
\end{table}

Атомарність складніше забезпечити у розподіленій системі, ніж у централізованій, оскільки одна логічна операція змінює дані на кількох незалежних вузлах. Один із вузлів може виконати локальну зміну, а інший --- втратити зв'язок або завершити оброблення з помилкою. Саме тому необхідний координатор та протокол, який унеможливлює фіксацію лише частини переказу.

Взаємне блокування може виникнути, якщо паралельні транзакції захоплюють ресурси в різному порядку: перша транзакція утримує рахунок~A й очікує рахунок~B, а друга утримує рахунок~B та очікує рахунок~A. У 2PC значення цього ризику зростає, оскільки prepared transaction продовжує утримувати блокування до отримання глобального рішення.

\paragraphheading{Критерії вибору механізму завершення транзакції.}
Вибір протоколу завершення розподіленої транзакції визначається не лише вимогою атомарності, а й допустимим часом блокування ресурсів, імовірністю відмов та можливістю відновити рішення координатора. Для операції переказу коштів пріоритетною є неможливість часткової фіксації: кінцевий стан має містити одночасно зменшення балансу відправника і збільшення балансу отримувача. Тому в моделі доцільно застосовувати 2PC, навіть якщо він передбачає очікування остаточного рішення після фази підготовки.

На фазі \texttt{PREPARE} кожний учасник перевіряє локальні обмеження: існування рахунку, його активність та достатність залишку на рахунку відправника. Після успішної перевірки вузол зберігає підготовлений стан у журналі та повідомляє координатору \texttt{READY}. З цього моменту локальна зміна ще не доступна як остаточно зафіксований результат, однак необхідні для неї ресурси можуть залишатися заблокованими до надходження команди завершення.

З практичної точки зору надійність 2PC потребує контролю за тривалістю prepared transactions. Координатор повинен зберігати глобальне рішення до його підтвердженого виконання кожним вузлом, а адміністратор --- мати можливість перевірити незавершені операції через системне подання \texttt{pg\_prepared\_xacts}. Таке поєднання протоколу, журналювання та діагностики забезпечує відновлюваність операції без порушення узгодженості даних.
